\section{Each Element in Its Setting}
\label{sec:each-element}

Before the interpretations, each text needs its own setting: where it stands in
its book, what its words are in the Latin and Greek behind the study
translation, and how the liturgical books relate it to the others. The books
state four relations. The first reading is chosen for its relation to the
Gospel, and the titles over the two show it (\work{Praenotanda} 106); the psalm
responds to the first reading (\work{General Instruction} 61); the Alleluia
verse greets the Gospel (\work{General Instruction} 62); and the second
reading and the Gospel each belong to a course read semi-continuously,
independent of the other. The Missal's antiphons and orations belong to the
Twenty-ninth Week in Ordinary Time and serve all three years of the
Lectionary's cycle.

Some words recur across these texts. The prayers ask to serve God and his gifts; the Gospel commands
rendering; the psalm bids the nations bring glory and sacrifices. The coin has
an image and an inscription, and in the Vulgate of Isaiah the Lord has made a
likeness of Cyrus and called him by name. Nations and their rulers appear in
the oracle to Cyrus, in the gods of the nations and the Lord who reigns among
them, in Caesar, in the crooked generation among whom the faithful shine, and,
in the verses around the two Communion antiphons, in the king not saved by his
army and the rulers of the gentiles. Cyrus has not known the Lord; the
questioners say ``we know''; Jesus knows their wickedness; the Lord acts so
that all may know. The Collect asks for a sincere heart, the Entrance
Antiphon's psalm speaks of lips without deceit, and the Gospel names
hypocrites.

\subsection{The Gospel: a question with two horns}

Matthew sets the scene in the temple, where the chief priests and the ancients
of the people had challenged Jesus' authority (21:23) and he had answered with
parables. After the second, the husbandmen of the vineyard, the chief priests
and Pharisees ``knew that he spoke of them'' and, seeking to lay hands on him,
feared the multitudes (21:45--46); the parable of the wedding feast followed
(22:1--14). The Pharisees now go aside and take counsel how to
catch him in his speech (22:15). They do not come themselves. They send their
disciples, and with them the Herodians, whom Jerome takes to be Herod's
soldiers or else those whom the Pharisees mocked for paying the tribute, and
whom Chrysostom calls ``Herod's soldiers''.

The question is built to wound on either answer. Aquinas sets out the two
opinions behind it: some held that a people dedicated to God ought not to be
tributary to a man, others that all owed Caesar tribute because he fought for
the peace of all; and so ``the question had two members: not to pay pertained
to the honour of God, to pay to the favour of men'' (\latin{Quaestio habebat
duo membra: scilicet quod non solverent, ad honorem Dei pertinebat; quod
solverent, ad favorem hominum}; \work{Super Matthaeum} c.~22, Venice ed.,
p.~283). Augustine calls it a ``double noose'': a yes would make Jesus a
maligner of the nation, a no a rebel against Caesar (\work{Enarrationes in
Psalmos} 57.8). Chrysostom says the plotters were ``digging, as they thought, a
precipice on either side'' (\work{Homilies on Matthew} 70.1). The praise that
opens the question serves the trap. They call him Master, truthful, a teacher
of God's way, one who regards no man's person; and in the last two phrases,
Chrysostom says, ``they were hinting at Herod and Caesar'', urging him toward
an answer that would set him against the rulers (70.1).

Jesus calls them hypocrites and asks for ``the coin of the tribute''. The
Greek names it \textit{to nomisma tou k\=ensou}, the coin of the census, and
they bring a denarius. He asks whose image and inscription it carries, and they
answer, ``Caesar's''. Jerome names the emperor: the Caesar of the saying is
Tiberius. The answer turns on a verb. The questioners had asked whether it is
lawful to \emph{give} (\latin{dare}) the census; Jesus commands them to
\emph{render} (\latin{reddite}, Greek \textit{apodote}), to give back what is
owed. Chrysostom builds on the change of verb (70.2), and the second
interpretation below follows him. Then comes the
second clause, ``and to God, the things that are God's'', which the questioners
had not asked about and which the saying does not explain. The reading ends
there. The next verse, in which the questioners wonder and leave him, is not
read.

The text of these verses is steady. The SBL Greek apparatus records only a
variant form of one participle and an omitted pronoun; the Clementine Vulgate
and the \work{Nova Vulgata} differ in a mood, a spelling and the naming of
Jesus in v.~20. Nothing in the
argument below depends on a variant.

\subsection{The first reading: an oracle to a foreign king}

Isaiah 45 stands among the chapters of the book that speak to Israel of its
release from Babylon (48:20). The verses just before the reading present the Lord as the
redeemer who made all things, who confounds diviners and confirms the word of
his servants, ``Who say to Cyrus: Thou art my shepherd, and thou shalt perform
all my pleasure. Who say to Jerusalem: Thou shalt be built: and to the temple:
Thy foundations shall be laid'' (44:24--28). Chapter 45 opens a new oracle
addressed to Cyrus himself. The Lord holds his right hand, subdues nations
before him and opens gates (v.~1). The verses the Lectionary leaves out promise
that the Lord will go before him, break the gates of brass and give him hidden
treasures, ``that thou mayest know that I am the Lord who call thee by thy
name, the God of Israel'' (vv.~2--3). The reading resumes with the reason: for
the sake of Jacob his servant the Lord has called Cyrus by name, though Cyrus
has not known him (v.~4); there is no God besides the Lord, who girded Cyrus
though Cyrus did not know him (v.~5); and the purpose is that all, from the
rising of the sun to its setting, may know that there is none besides him
(v.~6). The verses that follow, unread, go further: the Lord forms light and
creates darkness, makes peace and creates evil (v.~7); he has raised Cyrus up
``to justice'', and Cyrus ``shall build my city, and let go my captives, not
for ransom, nor for presents'' (v.~13).

Two words of v.~1 and v.~4 have a history. In v.~1 the Hebrew calls Cyrus the
Lord's anointed. The Clementine Vulgate renders the title with
\latin{christo}, a transliteration of the Greek, and the \work{Ordo}'s Latin
with \latin{uncto}; the sense is the same. An early reading, preserved in Latin
and Greek testimonies, took the Greek name of Cyrus (\textit{Kyr\=o}) as the
word for ``Lord'' (\textit{Kyri\=o}) and read the verse as said ``to Christ my
Lord''; the second interpretation below weighs it. In v.~4 the Clementine's
\latin{assimilavi te}, ``I have made a likeness of thee'', is not the
Hebrew's sense, which the King James Version gives as ``I have surnamed thee'';
Brenton's English of the Septuagint reads ``and accept thee''. The
Lectionary's English follows the Hebrew. A Father who reads \latin{assimilavi}
is reading the Vulgate's word.

The \work{Ordo}'s title for the reading, \latin{Apprehendi dexteram Cyri ut
subiciam ante faciem eius gentes}, condenses v.~1, putting the name of Cyrus
where the verse has a relative pronoun: what the Lectionary singles out is the
Lord's grasp of a gentile ruler's hand.

\subsection{The Responsorial Psalm: the nations brought to the Lord}

Psalm 96 (95) summons all the earth to a new song. It moves from the command to
sing and to declare the Lord's glory among the nations (vv.~1--3), through the
reason, that the Lord is great and to be feared above all gods, because the
gods of the nations are nothing and the Lord made the heavens (vv.~4--5), to
the summons of the families of the peoples to bring him glory and sacrifices
and to enter his courts (vv.~7--8), and to the proclamation among the nations
that the Lord reigns and will judge the peoples with justice (vv.~9--10). Its
close, not sung, calls heaven, earth, sea and trees to rejoice before the Lord
who comes to judge the world (vv.~11--13). Under the \work{General
Instruction} the psalm responds to the first reading; read beside it, it sets
the nations' recognition that their gods are nothing after the Lord's ``there
is none besides me'' (Is 45:5--6).

The Vulgate and Greek title, which the Hebrew psalm lacks, sets it ``when the
house was built after the captivity'', the rebuilding that Isaiah 44:28 and
45:13 promise through Cyrus. Augustine reads the house as the whole earth
built up of the nations: ``all the earth is the house of God'', and ``the very
act of singing is building'' (\work{Enarrationes in Psalmos} 95.1--2). In
v.~10 Augustine's psalter read ``the Lord hath reigned from the wood'', an Old
Latin addition that neither the Clementine nor the Lectionary's Latin reads,
so that what he says of it concerns his own text and not the appointed verse.
In v.~5 the Latin and Greek call the gods of the nations demons, as the
Douay--Rheims does, where the Hebrew calls them idols; Augustine and
Bellarmine expound the first form.

\subsection{The Second Reading: a church in God}

The narrative of Acts 17:1--10 sets the scene of the church Paul is
addressing. Paul had reasoned from the scriptures in the synagogue of
Thessalonica on three sabbath days; some believed and joined Paul and Silas,
``and of those that served God and of the Gentiles a great multitude''; a
tumult was raised, and Jason and certain brethren were drawn before the city's
rulers with the charge that ``these all do
contrary to the decrees of Caesar, saying that there is another king, Jesus''
(Acts 17:7), and Paul and Silas were sent away by night. Acts is not appointed
on this Sunday; it is the narrative behind the letter.

The reading is the letter's greeting and opening thanksgiving. Paul, with
Silvanus and Timothy, writes ``to the church of the Thessalonians: in God the
Father and in the Lord Jesus Christ''. Chrysostom notes that Paul gives himself
no title here because his readers were newly instructed, and he reads ``in
God'' as a demand: ``If any one walks not according to God, he cannot be said
to be `in God' '' (\work{Homilies on First Thessalonians}~1). Aquinas reads the
same phrase as the faith of the Trinity, \latin{in fide Trinitatis}. Paul
remembers their work of faith, labour of charity and enduring hope; the
\work{Ordo}'s title condenses this to \latin{Memores fidei vestrae et caritatis
et spei}. Aquinas says that faith ``is not sufficient unless it has work and
labour'' (\latin{non est sufficiens, nisi habeat operationem et laborem}) and
that this hope is ``before God, not before the eyes of men'' (\latin{Haec
spes est ante Deum, non ante oculos hominum}). Chrysostom draws
the labour of charity from the city's persecution: ``Merely to love is no labor
at all. But to love genuinely is great labor''. Paul knows their election
because the gospel came to them ``not \ldots\ in word only, but in power
also: and in the Holy Ghost and in much fulness''. Aquinas explains how he
knows: God gave him a great sign in the preaching, when he gives hearers the
grace to hear the word with fruit, and the preaching itself ``was not in
loftiness of speech, but in power'' (\latin{non fuit in sublimitate sermonis,
sed in virtute}; \work{Super I ad Thessalonicenses} c.~1, lect.~1).

The reading ends with the fullness of the gospel's coming. The clause that
follows, on what manner of men the missionaries were, is not read. The letter
is read on this Sunday because its course begins here; the Gospel's course
follows Matthew, and the two are independent.

\subsection{The Alleluia Verse: lights in a crooked generation}

The verse comes from Paul's exhortation in Philippians 2:12--18. The
Philippians are to work out their salvation with fear and trembling, since God
works in them both to will and to accomplish (vv.~12--13); to do all things
without murmuring and hesitation (v.~14), ``That you may be blameless and
sincere children of God, without reproof, in the midst of a crooked and
perverse generation: among whom you shine as lights in the world. Holding forth
the word of life'' (vv.~15--16). The Latin word behind ``sincere'' is
\latin{simplices}, single. The verse the Lectionary sets before the Gospel
takes only the shining and the word of life.

Chrysostom hears the image as a comparison with the stars: ``For the stars too
give light in the night, they shine in the dark, and receive no blemish to
their own beauty, yea they even shine the brighter \ldots\ Thus thou too dost
appear with the greater lustre, whilst thou holdest straight in the midst of
the crooked'' (\work{Homilies on Philippians}~8). Aquinas observes that
however the world changes, the heavenly lights stay bright, and that the
faithful are light not by their own essence, for only God is light in that
way, but ``inasmuch as they have something of the light of him who was the
light of men'' (\latin{sunt lux inquantum habent aliquid lucis illius qui erat
lux hominum}; \work{Super ad Philippenses} c.~2, lect.~4, Dessain ed.,
p.~383). Under the \work{General Instruction} the verse greets the Gospel; what
it says of the Gospel is drawn by the interpretations below.

\subsection{The Entrance Antiphon: the apple of the eye}

Psalm 17 (16), titled in the Latin ``The prayer of David'', is the plea of a
man who asks God to hear a just cause, ``my prayer, which proceedeth not from
deceitful lips'' (v.~1), and who says, ``Thou hast proved my heart, and
visited it by night'' (v.~3). The antiphon begins only at v.~6, with the cry
of one whom God has heard, and takes from v.~8 the request to be kept as the
apple of the eye and sheltered under God's wings. It passes over v.~7 and the
opening words of v.~8, which ask protection ``From them that resist thy right
hand''.

Augustine assigns the prayer ``to the Person of the Lord, with the addition of
the Church, which is His body'' (\work{Enarrationes in Psalmos} 16.1). On v.~6
he hears the confidence of one already answered: ``With a free and strong
effort have I directed My prayers unto Thee: for that I might have this power,
Thou hast heard Me when praying more weakly''. The apple of the eye, he says,
``seems very little and minute: yet by it is the sight of the eye directed,
whereby the light is distinguished from the darkness'', and the wings are
``the defence of Thy love and mercy'' (16.6, 8). Aquinas explains that the
pupil is guarded with diligence because nothing that could harm it is allowed
near, and reports the Gloss, which calls the pupil Christ, who discerns the
faithful from the unfaithful (\work{Postilla super Psalmos} 16). Bellarmine
calls the apple of the eye ``a most delicate, though valuable article'', and
those who pray ``frail and delicate'', needing the same care. The antiphon is a prayer of trust in
God's protection; its psalm's surrounding verses, unsung, concern sincerity of
lips and a heart that God has tried.

\subsection{The Collect}

The Collect is a single petition in two members, the will and the service. It
asks that the assembly always hold its will devoted to God, and that it serve
his majesty with a sincere heart. Its Latin stands word for word in the Old
Gelasian sacramentary as H.~A. Wilson edited it in 1894, where it is the first
prayer of the Mass headed \latin{Quinta Dominica post clausum Paschae}. It
recurs in the third book of the same sacramentary as the second prayer of an
\latin{Item alia Missa} (III.ix), word for word in Wilson's text; his
apparatus records that the Vatican manuscript reads \latin{maiestatem tuam}
there.

\subsection{The Prayer over the Offerings}

The Prayer over the Offerings asks that those who offer serve God's own gifts
with a free mind, and that his grace, purifying them, make them clean through
the mysteries they serve. The gifts are called God's (\latin{donis tuis})
before they are offered. The prayer stands in the Verona sacramentary as C.~L.
Feltoe edited it in 1896, among its Masses for April, with a single difference
of word order in the final clause (\latin{purificante nos gratia tua} for the
Missal's \latin{tua purificante nos gratia}).

\subsection{The Communion Antiphons}

The first antiphon takes two verses from the middle of Psalm 33 (32). Before
them, the Lord looks from heaven on all the sons of men, ``He who hath made the
hearts of every one of them: who understandeth all their works'' (vv.~13--15),
and ``The king is not saved by a great army: nor shall the giant be saved by
his own great strength. Vain is the horse for safety'' (vv.~16--17). Then: the
eyes of the Lord are on those who fear him and hope in his mercy, to deliver
them from death and feed them in famine (vv.~18--19). Augustine reads the king
not saved by strength as anyone ``who ruleth his own flesh, if he presume much
upon his own strength'', and the eyes of the Lord as his love: ``if thou seek
salvation, behold, the love of the Lord is upon them that fear Him''; the food
in famine is ``the nourishment of the Word, and of Everlasting Truth''
(\work{Enarrationes in Psalmos} 32.16, 18--19). Bellarmine joins the two
dispositions the verse names: ``For fear, without hope, is servile fear; hope,
without fear, is presumption. Fear, combined with hope, is the mark of real
love.'' Aquinas says the same in a sentence, \latin{timor sine spe desperat, et
spes sine timore praesumit} (\work{Postilla} 32).

The second antiphon is the last clause of Mark 10:45. Its context is the
request of James and John for the places at Jesus' right and left, the
indignation of the ten, and his answer: ``You know that they who seem to rule
over the Gentiles lord it over them: and their princes have power over them.
But it is not so among you'' (10:42--43); whoever would be first must be the
servant of all, ``For the Son of man also is not come to be ministered unto:
but to minister and to give his life a redemption for many'' (10:44--45). The
antiphon keeps only the giving of the life. Bede reads the passage as the
example of a humble and gentle teacher, and the verse of the one who ministered
and gave his life as a ransom, ``when he took the form of a servant, that he
might pour out his blood for the world'' (\latin{quando formam servi accepit,
ut pro mundo sanguinem funderet}); he restricts ``many'' to ``those who
willed to believe'' (\latin{pro eis qui credere voluerunt}; \work{In Marci
evangelium expositio} III, PL~92, col.~236). Chrysostom, on Matthew's form of
the same saying, says that the ransom was given ``for enemies''
(\work{Homilies on Matthew} 65.4). Bede speaks of those who receive the ransom
by faith, Chrysostom of what they were when it was paid.

The Missal's rubric for Ordinary Time lets either antiphon be chosen and
prefers one that agrees with the Gospel of the Mass. Neither is drawn from
Matthew 22:15--21, and the books do not say whether either agrees with it.

\subsection{The Prayer after Communion}

The Prayer after Communion asks that frequent sharing in heavenly things make
the assembly advance, aided by temporal benefits and instructed by eternal
ones. Its second half, from \latin{caelestium rerum frequentatione} to the
eternal things, stands in a prayer of the Verona sacramentary in the third
person, for a people that asks to rejoice in God's blessing; its opening words
are not there.
